\chapter{Supersymmetry Algebras and Representation Data}
\label{ch:susy-algebras-representations}

\ConventionsLorentzian

Supersymmetry is an extension of spacetime symmetry by odd generators.  The
word ``odd'' refers to a \(\mathbb Z_2\)-graded algebra: even generators
commute by ordinary commutators, odd generators anticommute with odd
generators, and mixed brackets are ordinary commutators.  This chapter fixes
the representation-theoretic data before introducing superfields or
Lagrangians.

\section{Particle Multiplets and Field-Variable Multiplets}

\begin{definition}[Particle supersymmetry representation]
\label{def:susy-particle-representation}
A particle representation of supersymmetry is a positive-energy unitary
representation of the real super-Poincare algebra on a \(\mathbb Z_2\)-graded
Hilbert space, or on a one-particle spectral subspace of such a Hilbert
space.  The odd generators are closed operators whose anticommutators are
understood on a common invariant domain.  A particle supermultiplet is an
irreducible representation after the mass orbit, spin or helicity data,
internal charges, and central charges have been fixed.
\end{definition}

\begin{definition}[Field-variable supersymmetry representation]
\label{def:susy-field-variable-representation}
A field-variable representation of supersymmetry is an action of the
super-Poincare algebra, or of its supertranslation subalgebra, by derivations
on a space of classical field variables, superfields, antifields, or local
functionals.  It is off shell when the algebra closes without imposing
Euler--Lagrange equations.  It is on shell when closure holds only modulo
equations of motion, gauge transformations, or both.  Auxiliary fields are
field variables added to obtain off-shell closure; they are not particle
states by that fact alone.
\end{definition}

The two notions of multiplet have different ambient categories.  Particle
multiplets live in Hilbert-space representation theory and are constrained by
positivity of the supercharge anticommutator.  Off-shell field multiplets live
in a coordinate description of local dynamics.  A relation between them is a
derived statement: one must choose an action or equations of motion, impose
constraints, quotient gauge redundancy if present, quantize or reconstruct
the state space, and then identify the resulting one-particle spectral
subspaces.  Counting component fields in an off-shell superfield is therefore
not a proof of the number or spin content of physical particles.

\section{Graded Lie Algebra}

\begin{definition}[Four-dimensional \(\mathcal N=1\) supertranslation algebra]
\label{def:n1-supertranslation-algebra}
Let \(P_\mu\) be the translation generators in the convention
\(U(a)=\exp(\ii a^\mu P_\mu)\).  The complexified \(\mathcal N=1\)
supertranslation algebra is generated by \(P_\mu\), odd Weyl spinors
\(Q_\alpha\), \(\bar Q_{\dot\alpha}\), and the identity, with
\[
  [P_\mu,P_\nu]=[P_\mu,Q_\alpha]=[P_\mu,\bar Q_{\dot\alpha}]=0,
\]
\[
  \{Q_\alpha,Q_\beta\}=0,
  \qquad
  \{\bar Q_{\dot\alpha},\bar Q_{\dot\beta}\}=0,
  \qquad
  \{Q_\alpha,\bar Q_{\dot\beta}\}
  =
  2\sigma^\mu_{\alpha\dot\beta}P_\mu .
\]
The matrices \(\sigma^\mu_{\alpha\dot\beta}\) and the raising/lowering
conventions are those of Section~\ref{sec:spinor-conventions}.  A real form is obtained by a
choice of adjoint operation for which \(P_\mu\) is Hermitian and
\((Q_\alpha)^\dagger=\bar Q_{\dot\alpha}\) after the stated spinor
identifications.
\end{definition}

The Poincare generators \(M_{\mu\nu}\) act on \(Q\) and \(\bar Q\) by the two
Weyl spinor representations.  This action is part of the full super-Poincare
algebra and fixes the Lorentz covariance of the anticommutator above.

\section{Positivity and Multiplet Size}

For a one-particle state of mass \(m>0\), go to its rest frame
\(P^\mu=(m,0,0,0)\).  The supercharge anticommutator becomes
\[
  \{Q_\alpha,Q_\beta^\dagger\}=2m\,\delta_{\alpha\beta}
\]
after choosing the rest-frame spinor basis.  Thus
\[
  a_\alpha=\frac{1}{\sqrt{2m}}Q_\alpha,\qquad
  a_\alpha^\dagger=\frac{1}{\sqrt{2m}}Q_\alpha^\dagger
\]
generate a two-mode fermionic oscillator algebra.

\begin{proposition}[Massive \(\mathcal N=1\) oscillator module]
\label{prop:massive-n1-oscillator-module}
Fix a massive one-particle orbit of mass \(m>0\), and fix an irreducible
spin-\(j\) representation \(V_j\) of the massive little group \(SU(2)\).  The
corresponding irreducible rest-frame supertranslation module is, as a
\(\mathbb Z_2\)-graded vector space,
\[
  V_j\otimes \Lambda^\bullet \mathbb C^2 .
\]
It has total dimension \(4(2j+1)\), bosonic dimension \(2(2j+1)\), and
fermionic dimension \(2(2j+1)\).  Under the diagonal little-group action the
fermionic one-oscillator sector carries
\[
  V_j\otimes \mathbb C^2
  \cong
  V_{j+1/2}\oplus V_{j-1/2},
  \label{eq:massive-n1-spin-decomposition}
\]
with the \(V_{j-1/2}\) summand omitted when \(j=0\), while the zero- and
two-oscillator sectors are two copies of \(V_j\) with opposite fermion
number parity from the one-oscillator sector.
\end{proposition}

\begin{proof}
The normalized operators \(a_\alpha,a_\alpha^\dagger\) satisfy
\[
  \{a_\alpha,a_\beta^\dagger\}=\delta_{\alpha\beta},
  \qquad
  \{a_\alpha,a_\beta\}=\{a_\alpha^\dagger,a_\beta^\dagger\}=0 .
\]
The irreducible representation of this Clifford algebra with a cyclic vector
annihilated by the \(a_\alpha\) is the exterior algebra
\(\Lambda^\bullet\mathbb C^2\).  Tensoring with the chosen little-group
representation gives the stated vector space and dimension.  The two
creation operators transform as a spin-\(1/2\) doublet, so the
one-oscillator sector is \(V_j\otimes V_{1/2}\), which decomposes by the
Clebsch--Gordan rule into \eqref{eq:massive-n1-spin-decomposition}.  The
zero-oscillator sector \(\Lambda^0\mathbb C^2\) and the two-oscillator sector
\(\Lambda^2\mathbb C^2\) are both little-group scalars, giving two copies of
\(V_j\).  Fermion parity is exterior degree modulo \(2\), hence the even and
odd dimensions are both \(2(2j+1)\).
\end{proof}

\begin{proposition}[Equal bosonic and fermionic dimensions]
\label{prop:susy-equal-boson-fermion}
In any finite-dimensional representation of the massive rest-frame
\(\mathcal N=1\) oscillator algebra with positive inner product, the bosonic
and fermionic state-space dimensions are equal.
\end{proposition}

\begin{proof}
Let \(F=(-1)^{N_F}\) be the fermion-parity operator of the oscillator algebra.
For a two-mode fermionic Fock space,
\[
  \operatorname{Tr}F=(1-1)^2=0 .
\]
Tensoring with a little-group representation multiplies the trace by its
dimension, so the trace remains zero.  Since \(F\) has eigenvalue \(+1\) on
bosonic states and \(-1\) on fermionic states, \(\operatorname{Tr}F=0\) is
exactly equality of the two dimensions.
\end{proof}

\section{Massless Representations}

For a massless momentum \(p^\mu=(E,0,0,E)\), the matrix
\(\sigma^\mu p_\mu\) has rank one.  Only one complex linear combination of
the supercharges acts with nonzero norm on a positive-energy representation;
the null combination is quotiented.  The remaining creation operator changes
helicity by \(1/2\).  Hence a massless irreducible representation is built
from a Clifford algebra with one fermionic oscillator, and its helicities
come in adjacent pairs \((h,h+1/2)\) before CPT completion.

\begin{proposition}[Massless \(\mathcal N=1\) helicity pair]
\label{prop:massless-n1-helicity-pair}
Let \(E>0\) and \(p^\mu=(E,0,0,E)\).  In a spinor basis where
\(\sigma^\mu p_\mu\) is diagonal, positivity sets one complex supercharge
combination to zero on the Hilbert space and leaves one normalized
fermionic oscillator.  An irreducible massless \(\mathcal N=1\)
representation generated from a helicity-\(h\) state therefore has two
helicity states,
\[
  h,\qquad h+\frac12,
\]
before imposing CPT.
\end{proposition}

\begin{proof}
For the chosen momentum, \(\sigma^\mu p_\mu\) is a positive semidefinite
\(2\times2\) Hermitian matrix of rank one.  A unitary change of spinor basis
therefore puts the supercharge norm matrix into the form
\[
  \{Q_a,Q_b^\dagger\}
  =
  \begin{pmatrix}
    cE&0\\
    0&0
  \end{pmatrix}_{ab},
  \qquad c>0 .
\]
If \(\psi\) is any vector in the common invariant domain, then
\[
  0
  =
  \langle \psi,\{Q_2,Q_2^\dagger\}\psi\rangle
  =
  \|Q_2\psi\|^2+\|Q_2^\dagger\psi\|^2
\]
because the second diagonal entry is zero.  Positivity of the inner product
therefore forces \(Q_2\psi=Q_2^\dagger\psi=0\) for all vectors in the
domain, equivalently the null supercharge combination acts trivially after
quotienting zero-norm vectors.  The remaining normalized operator
\[
  b=(cE)^{-1/2}Q_1,\qquad b^\dagger=(cE)^{-1/2}Q_1^\dagger
\]
satisfies \(\{b,b^\dagger\}=1\).  Because a Weyl supercharge has helicity
\(\pm1/2\), the nontrivial creation operator shifts helicity by one half in
the chosen convention.  Thus the irreducible Clifford module has precisely
the two displayed helicities.
\end{proof}

\section{Central Charges and BPS Bounds}

For extended supersymmetry, odd generators carry an index \(I=1,\ldots,N\).
The algebra may contain antisymmetric central charges,
\[
  \{Q^I_\alpha,Q^J_\beta\}
  =
  \epsilon_{\alpha\beta}Z^{IJ},
  \qquad Z^{IJ}=-Z^{JI}.
\]
In a rest frame, positivity of the matrix of anticommutators gives lower
bounds on the mass in terms of the singular values of \(Z\).  A state
saturating such a bound is annihilated by a nonzero linear combination of
supercharges; its multiplet is shorter because the corresponding fermionic
oscillator is absent in the quotient representation.

\begin{proposition}[Rest-frame BPS bound from singular values]
\label{prop:bps-bound-singular-values}
In a massive rest frame for extended supersymmetry, suppose the central-charge
matrix has singular values \(z_s\ge0\) after skew-diagonalization of the
antisymmetric matrix \(Z^{IJ}\).  Positivity of the Hilbert-space
representation implies
\[
  m\ge z_s
  \qquad\text{for every singular value }z_s .
\]
If \(m=z_s\) for some \(s\), then a nonzero linear combination of
supercharges has zero norm and annihilates the shortened representation.
\end{proposition}

\begin{proof}
The only linear-algebra input is the skew singular-value decomposition for a
complex antisymmetric matrix: a unitary change of the extended-supersymmetry
index brings \(Z^{IJ}\) to a direct sum of \(2\times2\) blocks
\(z_s\varepsilon\), \(z_s\ge0\), plus possible zero blocks.  In the associated
real charge combinations the positive anticommutator matrix decomposes into
independent Hermitian blocks of the form
\[
  2
  \begin{pmatrix}
    m I_r & Z_s\\
    Z_s^\dagger & m I_r
  \end{pmatrix},
  \qquad
  Z_s=z_s I_r
\]
on the subspace associated with the singular value \(z_s\), where \(I_r\) is
the identity on its multiplicity space.  The eigenvectors of the displayed
Hermitian block are the symmetric and antisymmetric combinations of the two
entries, and the corresponding eigenvalues are
\[
  2(m+z_s),\qquad 2(m-z_s).
\]
A positive-energy unitary representation requires every anticommutator norm
matrix to be positive semidefinite, hence \(m-z_s\ge0\).  If \(m=z_s\), the
second eigenvalue is zero.  The associated supercharge combination has zero
norm on every vector in the representation; after quotienting zero-norm
vectors it acts by zero.  Removing that fermionic oscillator is precisely the
shortening mechanism.
\end{proof}

\begin{openproblem}[Nonperturbative realization of supersymmetric QFT data]
\label{op:susy-nonperturbative-realization}
For a proposed supersymmetric QFT, one must construct the Hilbert space or
Euclidean/functorial object, the supercharges as operators or Ward
identities, the stress tensor and supercurrent, and the required positivity
structure.  Perturbative superfield expressions and protected-index
calculations are powerful coordinates on such data, but a nonperturbative
construction requires separate control of the regulator, local operator
domains, anomalies, and possible supersymmetry-breaking effects.
\end{openproblem}
